\documentclass[10pt,a4paper]{amsart}
\usepackage[margin=19mm]{geometry}
\usepackage{amsmath,amssymb,mathtools}
\usepackage[scr=rsfs,cal=euler]{mathalfa}
\usepackage{xfrac}
\usepackage[unicode,hidelinks,bookmarks=false]{hyperref}
\newcommand{\ie}{i.e.\ }
\newcommand{\cf}{cf.\ }
\newcommand{\resp}{respectively\ }
\newcommand{\Subsection}{\subsection}
\providecommand{\nobreakdash}{\nobreak\hbox{-}}
\setlength{\emergencystretch}{2em}
\multlinegap0pt
\usepackage{mathtools}
\usepackage{stmaryrd}
\usepackage{enumitem}
\usepackage{tikz}
\usetikzlibrary{positioning,arrows.meta,calc}
\newtheorem{definition}{Definition}[section]
\newtheorem{example}[definition]{Example}
\newtheorem{remark}[definition]{Remark}
\newtheorem{fact}[definition]{Fact}
\newtheorem{theorem}[definition]{Theorem}
\newtheorem{claim}[definition]{Claim}
\newtheorem{lemma}[definition]{Lemma}
\newtheorem{proposition}[definition]{Proposition}
\newtheorem{corollary}[definition]{Corollary}
\newcommand{\Pow}{\mathcal P}
\newcommand{\Prop}{\mathsf{Prop}}
\newcommand{\Sent}{\mathsf{Sent}}
\newcommand{\SentP}{\mathsf{Sent}_P}
\newcommand{\LP}{L_{\in}} % language of set theory
\newcommand{\LB}{L_B}     % parameter language
\newcommand{\Lbox}{\mathcal L_{\Box}} % modal language
\newcommand{\val}[1]{\llbracket #1\rrbracket_B}
\newcommand{\Btwo}{\mathbf{B}_2}
\newcommand{\Bfour}{\mathbf{B}_4}
\newcommand{\KB}{\mathsf{KB}}
\newcommand{\KTB}{\mathsf{KTB}}
\newcommand{\Logfull}{\mathsf{Log}_{\mathrm{full}}^{\sharp}}
\newcommand{\Logsub}{\mathsf{Log}_{\subseteq}^{\sharp}}
\begin{document}
\title[The Internal Modal Logic of Forcing]{The Internal Modal Logic of Forcing}
\author{Santiago Jockwich, Sourav Tarafder; Giorgio Venturi}
\date{}
\maketitle
\noindent\emph{Source and licence.} The Internal Modal Logic of Forcing. Version arXiv:2607.25977v1 (CC BY 4.0). This material is distributed under the Creative Commons Attribution 4.0 International licence, http://creativecommons.org/licenses/by/4.0/, as recorded in the arXiv OAI licence metadata for this exact version and shown on the abstract page https://arxiv.org/abs/2607.25977v1. Full rights notice: licenses/arxiv-2607.25977-CC-BY-4.0.txt.\par\smallskip
\noindent\textit{Editorial scope.} Counted unit: the original \texttt{theorem} environment \texttt{thm:general-valid} at source lines 463--486 together with its complete proof at lines 488--520; the delivered document keeps source lines 142--762 so that every referenced label and every citation resolves inside it. Native editable source is the paper's own concatenated TeX; the AMS wrapper is editorial formatting only. No publisher class, upstream script or unrelated artwork is redistributed.
\section{Technical preliminaries}

We briefly recall the algebraic notions used throughout.

\begin{definition}
A \emph{lattice} is an algebra $\langle A,\wedge,\vee\rangle$ where $A$ is nonempty and
$\wedge,\vee$ are binary operations satisfying commutativity, associativity, idempotence, and
absorption.
\end{definition}

A lattice induces a partial order by $a\le b$ iff $a\wedge b=a$.

\begin{definition}
A lattice $\langle A,\wedge,\vee\rangle$ is \emph{bounded} if it has a greatest element $1$ and
a least element $0$. It is \emph{complete} if every subset $S\subseteq A$ has a supremum $\bigvee S$ and an infimum $\bigwedge S$.
\end{definition}

\begin{definition}
A lattice $\langle A,\wedge,\vee\rangle$ is \emph{distributive} if for any $a, b, c \in A$,
\[a \wedge (b \vee c) = (a \wedge b) \vee (a \wedge c) \mbox{ and } a \vee (b \wedge c) = (a \vee b) \wedge (a \vee c).\]
\end{definition}

\begin{definition}
A \emph{Boolean algebra} is a bounded distributive lattice
$\langle A,\wedge,\vee,{}^\ast,1,0\rangle$ in which every $a\in A$ has a complement $a^\ast$
satisfying $a\wedge a^\ast=0$ and $a\vee a^\ast=1$. The following two abbreviations are used in a Boolean algebra:
\[a \to b = a^\ast \vee b \mbox{ and } a \leftrightarrow b = (a \to b) \wedge (b \to a).\]
\end{definition}

\begin{example}
The Boolean algebra with universe $\{0,1\}$ is denoted by $\Btwo$ and corresponds to classical
two-valued semantics.
\end{example}

\begin{example}
For any set $X$, the power set $\Pow(X)$ with $\cap,\cup,{}^c,X,\varnothing$ is a complete Boolean
algebra.
\end{example}

\begin{definition}
An element $a\ne 0$ is an \emph{atom} if there is no $b$ with $0<b<a$. A Boolean algebra $B$ is \emph{atomic} if for every element $a \in B$ there exists a subset $S \subseteq B$ of atoms such that $\bigvee S = a$.
\end{definition}

\begin{definition}\label{def:filter}
Let $B=\langle A,\wedge,\vee,1,0\rangle$ be a Boolean algebra.
A set $F\subseteq A$ is a \emph{filter} if:
\begin{enumerate}[label=(\roman*)]
\item $1\in F$ and $0\notin F$;
\item if $x\in F$ and $x\le y$, then $y\in F$;
\item if $x,y\in F$, then $x\wedge y\in F$.
\end{enumerate}
A filter $U$ is an \emph{ultrafilter} if for every $a\in A$, either $a\in U$ or $a^\ast\in U$.
\end{definition}

\subsection{Boolean-valued models}

Boolean-valued models were introduced by Vop\v{e}nka, Solovay, and Scott as a semantic companion
to forcing. We follow standard references \cite{Bell,Jech}.

\begin{definition}
Let $B$ be a complete Boolean algebra. Define by transfinite recursion:
\[
V^{(B)}_\alpha=\{x:\ x \text{ is a function, }\mathrm{ran}(x)\subseteq B,\ \exists\xi<\alpha\ (\mathrm{dom}(x)\subseteq V^{(B)}_\xi)\},
\]
and put $V^{(B)}=\bigcup_\alpha V^{(B)}_\alpha$.
Let $\LB$ be the expansion of $\LP$ obtained by adding constant symbols for each element of $V^{(B)}$. And let $\Sent(L_B)$ be the class of corresponding sentences.
\end{definition}

\begin{definition}\label{def:boolean-interpretation}
The Boolean interpretation map $\val{\cdot}:\Sent(L_B)\to B$ is defined by recursion on formulas, starting
from membership and equality for $u,v\in V^{(B)}$:
\[
\begin{aligned}
\val{u\in v} \ &=\ \bigvee_{x\in\mathrm{dom}(v)}\bigl(v(x)\wedge \val{x=u}\bigr),\\[2pt]
\val{u=v} \ &=\ \Bigl(\bigwedge_{x\in\mathrm{dom}(u)}\bigl(u(x)\to \val{x\in v}\bigr)\Bigr)\ \wedge\
            \Bigl(\bigwedge_{y\in\mathrm{dom}(v)}\bigl(v(y)\to \val{y\in u}\bigr)\Bigr).
\end{aligned}
\]
and extending to the Boolean connectives and quantifiers by:
\[
\val{\varphi\wedge\psi}=\val{\varphi}\wedge\val{\psi},\quad
\val{\varphi\vee\psi}=\val{\varphi}\vee\val{\psi},\quad
\val{\neg\varphi}=\val{\varphi}^\ast,
\]
\[
\val{\forall x\,\varphi(x)}=\bigwedge_{u\in V^{(B)}}\val{\varphi(u)},
\qquad
\val{\exists x\,\varphi(x)}=\bigvee_{u\in V^{(B)}}\val{\varphi(u)}.
\]
\end{definition}

\begin{definition}
Write $V^{(B)}$ for the structure together with the interpretation map $\val{\cdot}$.
Given a filter $F$ on $B$, an $\LB$-sentence $\sigma$ is \emph{valid in $V^{(B)}$ relative to $F$} if
$\val{\sigma}\in F$, written $V^{(B)}\models_F \sigma$.
\end{definition}

\begin{theorem}[\cite{Bell}]\label{thm:zfc-boolean}
For every complete Boolean algebra $B$, the Boolean value of each axiom of $\mathsf{ZFC}$ is $1$.
Equivalently, $V^{(B)}\models \mathsf{ZFC}$ (in the sense $\val{\mathsf{ZFC}}=1$).
\end{theorem}

To show that a sentence in the language of set theory $\sigma$ is consistent with $\mathsf{ZFC}$ using Boolean-valued
models, it suffices to exhibit a complete $B$ with $\val{\sigma}\ne 0$.
Indeed, the principal filter generated by $\val{\sigma}$ extends to an ultrafilter $U$ on $B$ with
$\val{\sigma}\in U$, and the corresponding quotient $V^{(B)}/U$ is a classical model of $\mathsf{ZFC}+\sigma$
(see Theorem~\ref{thm:ultrafilter-truth-lemma}).

\subsection{Faithfulness and loyalty}

Following \cite{LowePassmannTarafder}, we recall two notions measuring how much of the Boolean algebra $B$
is realized as truth values of \emph{parameter-free} sentences.

\begin{definition}[\cite{LowePassmannTarafder}]\label{def:faithful-loyal}
Let $B$ be a complete Boolean algebra and $F$ a filter on $B$.
The Boolean-valued model $V^{(B)}$ is \emph{loyal} to $(B,F)$ if the propositional logic of $(B,F)$
coincides with the propositional logic induced by $V^{(B)}$ relative to $F$.
It is \emph{faithful to $B$} if for every $a\in B$ there exists a sentence in the language of set theory $\sigma$ such that $\val{\sigma}=a$.
\end{definition}

\begin{lemma}[\cite{LowePassmannTarafder}]\label{lem:faithful-implies-loyal}
If $V^{(B)}$ is faithful to $B$ then it is loyal to $(B,F)$ for every filter $F$ on $B$.
\end{lemma}

\begin{theorem}[\cite{LowePassmannTarafder}]\label{thm:atomic-not-faithful}
If $B$ is an atomic complete Boolean algebra with more than two elements, then $V^{(B)}$ is loyal but not
faithful for parameter-free sentences. In particular, for every $\sigma$, sentence in the language of set theory,  we have
$\val{\sigma}\in\{0,1\}$.
\end{theorem}

\subsection{Modal logic}

We use standard propositional modal logic and stratify our language as usual.
The modal language is $\Lbox$ with propositional variables belonging to a countable set that we name $\Prop$ and Boolean connectives, and $\Box,\Diamond$.

\begin{definition}
A \emph{Kripke frame} is a pair $\langle W,R\rangle$ where $W$ is nonempty and $R\subseteq W\times W$.
A \emph{Kripke model} is a triple $M=\langle W,R,V\rangle$, where $V:\Prop\to \Pow(W)$ is a valuation.
\end{definition}

\begin{definition}\label{def:sat}
Given $M=\langle W,R,V\rangle$ and $w\in W$, define $M,w\models\varphi$ by the usual clauses:
\[
\begin{aligned}
&M,w\models p \iff w\in V(p),\\
&M,w\models \Box\varphi \iff \forall v\in W\ (wRv\Rightarrow M,v\models\varphi),\\
&M,w\models \Diamond\varphi \iff \exists v\in W\ (wRv\ \&\ M,v\models\varphi),
\end{aligned}
\]
and the standard Boolean clauses for $\neg,\wedge,\vee,\to$.
\end{definition}

\begin{definition}
A formula $\varphi$ is \emph{valid in a model} $M$, written $M\models \varphi$, if $M,w\models\varphi$ for all
$w\in W$. It is \emph{valid in a frame} $\langle W,R\rangle$ if it is valid in every model based on that frame.
\end{definition}

\section{Extended Boolean-valued semantics}\label{sec:extended-semantics}

In order to interpret the modal language in set theory we will interpret the \emph{atomic} modal propositions as (parameter-free) set-theoretic sentences. Let us call such a collection $\Sent_{\in}$.
Hence, fix an enumeration of $\Sent_{\in}$ as $\langle p_i:i\in\omega\rangle$, which covers all elements of $\Prop$.\footnote{For completeness, we later quantify over \emph{all} such enumerations; see Section~\ref{sec:completeness}.}

\begin{remark}
Notice that we are here using the symbols for connectives ($\neg, \land,  \lor, \to$) to indicate three different operations: 1) the algebraic operations, 2) the connectives in the language $L_\Box$, and 3) the connectives in the language $L_\in$. We are confident that the difference will be clear from the context.
\end{remark}

\subsection{Models, satisfaction, and theories}

We now need to introduce the relevant definition for interpreting our modal language.

\begin{definition}\label{def:MB}
Let $B$ be a complete Boolean algebra and let $R$ be an accessibility relation on $B$.
Define $M_B=\langle B,R,\val{\cdot}\rangle$ and interpret modal formulas at states $b\in B$ by:
\begin{enumerate}[label=(\roman*)]
\item $M_B,b\models p$ iff $b\le \val{p}$ for $p\in\Sent_{\in}$;
\item Boolean connectives as in Definition~\ref{def:sat};
\item $M_B,b\models \Box\varphi$ iff for all $a\in B$ with $bRa$, we have $M_B,a\models \varphi$;
\item $M_B,b\models \Diamond\varphi$ iff there exists $a\in B$ with $bRa$ and $M_B,a\models\varphi$.
\end{enumerate}
\end{definition}

\begin{definition}\label{def:theories}
Let $B$ be complete.
A modal formula $\varphi$ is \emph{valid in $M_B$} if $M_B,b\models\varphi$ for all $b\in B$; write $M_B\models\varphi$.
Define
\[
T_H(M_B)=\{\varphi\in\Lbox:\ M_B\models\varphi\}.
\]
We also consider \emph{validity}$^\sharp$ excluding the isolated state $0$: write $M_B\models^\sharp\varphi$ if
$M_B,b\models\varphi$ for all $b\in B\setminus\{0\}$, and define
\[
T_H^\sharp(M_B)=\{\varphi\in\Lbox:\ M_B\models^\sharp\varphi\}.
\]
\end{definition}

\begin{fact}\label{fact:zfc-sentences}
If $\sigma\in\Sent_{\in}$ has Boolean value $\val{\sigma}=1$, then $M_B\models \sigma$.
\end{fact}

\begin{proof}
If $\val{\sigma}=1$, then for every $b\in B$ we have $b\le 1=\val{\sigma}$, hence $M_B,b\models\sigma$ by
Definition~\ref{def:MB}(i).
\end{proof}

\subsection{Accessibility as co-consistency}\label{subsec:coconsistency}

We now introduce the accessibility relation that captures our intended internal notion of possibility.

\begin{definition}\label{def:R}
Let $B$ be a complete Boolean algebra. Define $R$ on $B$ by
\[
aRb \quad\Longleftrightarrow\quad a\wedge b\neq 0.
\]
\end{definition}

\begin{remark}\label{rem:R-equivalences}
Let $B$ be a Boolean algebra and $a,b\in B$. The following are equivalent:
\begin{enumerate}[label=(\roman*)]
\item $aRb$, i.e.\ $a\wedge b\neq 0$;
\item $b\not\le a^{\ast}$;
\item there exists an ultrafilter $U$ on $B$ such that $a, b\in U$.
\end{enumerate}
\end{remark}

\begin{proof}
(i)$\Leftrightarrow$(ii) is standard: $b\le a^\ast$ iff $a\wedge b=0$.
For (i)$\Rightarrow$(iii), extend the proper filter generated by $a\wedge b$ to an ultrafilter.
For (iii)$\Rightarrow$(i), if $a,b\in U$ then $a\wedge b\in U$, hence $a\wedge b\neq 0$.
\end{proof}

\begin{example}
Let $B=\Bfour=\{0,a,b,1\}$ with $a^\ast=b$ and $b^\ast=a$.
Then $0$ is $R$-isolated, $1$ is $R$-related to every nonzero element, and $aRb$ \& $bRa$ fail, since $a\wedge b=0$.
\end{example}

Notice that $0$ is an isolated state.

\begin{lemma}\label{lem:state0}
Let $B$ be a complete Boolean algebra. Then for every modal formula $\varphi\in\Lbox$:
\begin{enumerate}[label=(\roman*)]
\item $M_B,0\models \Box\varphi$;
\item $M_B,0\not\models \Diamond\varphi$.
\end{enumerate}
Moreover, for every atomic $p\in\Sent_{\in}$, we have $M_B,0\models p$.
\end{lemma}

\begin{proof}
There is no $a\in B$ with $0Ra$ since $0\wedge a=0$, so $\Box\varphi$ holds vacuously at $0$ and $\Diamond\varphi$
fails. For the atomic clause, $0\le \val{p}$ always holds.
\end{proof}

On the contrary $1$ is connected with every other state.

\begin{proposition}\label{prop:1-related}
If $a\in B\setminus\{0\}$, then $1Ra$ and $aR1$.
\end{proposition}

\begin{proof}
If $a\neq 0$, then $1\wedge a=a\neq 0$, so $1Ra$; and $a\wedge 1=a\neq 0$, so $aR1$.
\end{proof}

We now make precise the connection between internal possibility and the usual ultrafilter collapse.

\begin{definition}\label{def:ultrafilter-quotient}
Let $B$ be a complete Boolean algebra and $U$ an ultrafilter on $B$.
Define an equivalence (class) relation $\equiv_U$ on $V^{(B)}$ by
\[
u\equiv_U v \quad\Longleftrightarrow\quad \val{u=v}\in U,
\]
and define a membership relation $\in_U$ on $\equiv_U$-classes by
\[
[u]_U\in_U [v]_U \quad\Longleftrightarrow\quad \val{u\in v}\in U.
\]
Write $V^{(B)}/U$ for the resulting two-valued structure.
\end{definition}

Thus truth in a structure can be defined in terms of ultrafilters.

\begin{theorem}\label{thm:ultrafilter-truth-lemma}
For every $\LB$-sentence $\sigma$ and every ultrafilter $U$ on $B$,
\[
V^{(B)}/U\models \sigma \quad\Longleftrightarrow\quad \val{\sigma}\in U.
\]
\end{theorem}

\begin{proof}
This is the standard Boolean-valued \L o\'s lemma for ultrafilter quotients, proved by induction on the
complexity of formulas; see \cite{Bell,Jech}.
\end{proof}

And this reading provides a direct interpretation of the $\Diamond$-modality when applied to set-theoretical sentences.

\begin{proposition}\label{prop:diamond-ultrafilter-reading}
Let $p\in\Sent_{\in}$ (viewed as an atomic modal proposition) and let $b\in B$.
Then
\[
M_B,b\models \Diamond p
\quad\Longleftrightarrow\quad
\exists\,U\text{ an ultrafilter on }B\ \bigl(b\in U\ \wedge\ V^{(B)}/U\models p\bigr).
\]
\end{proposition}

\begin{proof}
($\Rightarrow$) Assume $M_B,b\models \Diamond p$. Then there is $a\in B$ with $bRa$ and $M_B,a\models p$.
By Definition~\ref{def:MB}(i), $M_B,a\models p$ means $a\le \val{p}$.
By Remark~\ref{rem:R-equivalences}(iii), choose an ultrafilter $U$ with $a,b\in U$.
Since $U$ is upward closed and $a\le \val{p}$, we have $\val{p}\in U$.
By Theorem~\ref{thm:ultrafilter-truth-lemma}, $V^{(B)}/U\models p$.

($\Leftarrow$) Conversely, suppose $U$ is an ultrafilter with $b\in U$ and $V^{(B)}/U\models p$.
By Theorem~\ref{thm:ultrafilter-truth-lemma}, $\val{p}\in U$.
Let $a:=\val{p}$. Then $a\in U$ and $b\in U$, hence $a\wedge b\in U$ and therefore $a\wedge b\neq 0$.
Thus $bRa$ by Definition~\ref{def:R}.
Also $M_B,a\models p$ holds because $a\le \val{p}$.
Therefore $M_B,b\models \Diamond p$.
\end{proof}

\section{Preliminary results}\label{sec:results}

Several modal principles hold for every complete Boolean algebra.


\begin{theorem}\label{thm:general-valid}
Let $B$ be a complete Boolean algebra. Then the following modal principles are in $T_H(M_B)$:
\[
\Box(\varphi\to\psi)\to (\Box\varphi\to \Box\psi)\qquad (K),
\]
\[
\neg\Diamond\varphi \leftrightarrow \Box\neg\varphi \qquad (\mathrm{Dual}),
\]
\[
\varphi \to \Box\Diamond\varphi \qquad (B),
\]
\[
\Diamond\Box\varphi \to \Box\Diamond\varphi\qquad (.2),
\]
and the ``symmetric'' principle
\[
\Diamond\Box\varphi \to \varphi \qquad (5^\ast).
\]
Consequently, the derived principle
\[
\Box\bigl(\Box(\varphi\to\Box\varphi)\to\varphi\bigr)\to(\Diamond\Box\varphi\to\varphi)\qquad (DM),
\]
also belongs to $T_H(M_B)$.
\end{theorem}

\begin{proof}
(K) and (Dual) are valid in all Kripke models.

For (B), assume $M_B,a\models \varphi$ and let $b$ be any state with $aRb$.
By symmetry of $R$ (since $a\wedge b\neq 0$ iff $b\wedge a\neq 0$), we have $bRa$.
Thus $M_B,b\models \Diamond\varphi$, so $M_B,a\models \Box\Diamond\varphi$.

For $(.2)$, assume $M_B,a\models \Diamond\Box\varphi$.
Choose $b$ with $aRb$ and $M_B,b\models \Box\varphi$.
Let $c$ be any state with $aRc$.
Put $d:=b\vee c$.
Since $aRc$, we have $c\neq 0$, and since $aRb$, we have $b\neq 0$.
Hence
\[
b\wedge d=b\neq 0
\qquad\text{and}\qquad
c\wedge d=c\neq 0,
\]
so $bRd$ and $cRd$.
Because $M_B,b\models \Box\varphi$ and $bRd$, we get $M_B,d\models \varphi$.
Thus $M_B,c\models \Diamond\varphi$.
As $c$ was arbitrary, $M_B,a\models \Box\Diamond\varphi$.



For $(5^\ast)$, assume $M_B,a\models \Diamond\Box\varphi$.
Then there exists $b$ with $aRb$ and $M_B,b\models \Box\varphi$.
By symmetry, $bRa$, so $M_B,b\models \Box\varphi$ implies $M_B,a\models \varphi$.



Finally, $(DM)$ is an immediate consequence since its consequent is $(5^\ast)$.
\end{proof}


On the other hand, the presence of the isolated state $0$ leads to immediate failures of some familiar principles.

\begin{theorem}\label{thm:general-fail}
Let $B$ be a complete Boolean algebra. Then the following modal principles are not in $T_H(M_B)$:
\[
\Box\varphi \to \varphi \qquad (T),
\qquad
\Box\varphi\to\Diamond\varphi \qquad (M),
\]
\[
\quad
\Box\bigl(\Box(\varphi\to \Box\varphi)\to \varphi\bigr)\to \varphi \qquad (\mathrm{Grz}).
%\qquad
%\Diamond\Box\varphi \to \Box\Diamond\varphi\qquad (.2).
\]
%and also the converse of $(.2)$ may fail as well.
\end{theorem}

\begin{proof}
Let $p\in\Sent_{\in}$ be any sentence and set $\varphi:=\neg p$.
By Lemma~\ref{lem:state0}, $M_B,0\models \Box\varphi$ but $M_B,0\not\models \varphi$.
Hence $(T)$ fails at $0$.
Also $M_B,0\not\models \Diamond\varphi$ by Lemma~\ref{lem:state0}$\,$(ii), so $(M)$ fails at $0$ as well.
For (Grz) the failure can be witnessed similarly by using the dead-end behavior of $0$.
\end{proof}

\begin{remark}\label{rem:nonzero-bridge}
If one works with $T_H^\sharp(M_B)$ (validity on $B\setminus\{0\}$), then $(T)$ and $(M)$ both hold because
$R$ is reflexive on $B\setminus\{0\}$ (indeed $aRa$ for all $a\ne 0$). This nonzero-state perspective
underlies the translation-based completeness theorem of Section~\ref{sec:completeness}.
\end{remark}

At this stage one should notice that different Boolean algebras can give different theories.

\begin{theorem}\label{thm:different-theories}
There exist complete Boolean algebras $B$ and $B'$ such that $T_H(M_B)\ne T_H(M_{B'})$.
\end{theorem}

\begin{proof}
Let $B=\Btwo$. In $M_{\Btwo}$, the frame has only two states, and $1$ accesses only itself.
One checks that the formula $p\to \Box p$ is valid in $M_{\Btwo}$.

Now let $B'$ be a complete Boolean algebra and choose a set-theoretic sentence $p\in\Sent_{\in}$ such that
$0<\val{p}<1$ in $B'$. (For example, choose $B'$ arising from a forcing for which $p$ is independent.)
Let $a:=\val{p}$.
Then $M_{B'},a\models p$ but $M_{B'},1\not\models p$, and since $aR1$ we have $M_{B'},a\not\models \Box p$.
Hence $p\to\Box p$ fails in $M_{B'}$, so $T_H(M_B)\ne T_H(M_{B'})$.
\end{proof}

\subsection{Complete atomic Boolean algebras}

When $B$ is complete and atomic, the restriction that atomic propositions are \emph{parameter-free}
set-theoretic sentences has a strong consequence: by Theorem~\ref{thm:atomic-not-faithful} every such
sentence has Boolean value $0$ or $1$. This collapses the atomic behavior and yields extra modal validities.

\begin{theorem}\label{thm:atomic-uniformity}
Let $B$ be a complete atomic Boolean algebra and let $\varphi\in\Lbox$.
Then either $M_B,a\models \varphi$ for every $a\in B\setminus\{0,1\}$, or $M_B,a\not\models \varphi$ for every
$a\in B\setminus\{0,1\}$.
\end{theorem}

\begin{proof}
We proceed by induction on the complexity of $\varphi$.

\smallskip\noindent\emph{Base case.}
If $\varphi$ is atomic, say $\varphi=p\in\Sent_{\in}$, then by Theorem~\ref{thm:atomic-not-faithful} we have
$\val{p}\in\{0,1\}$. If $\val{p}=0$ then no $a\ne 0$ satisfies $p$; if $\val{p}=1$ then every $a$ satisfies $p$.
In either case, all $a\in B\setminus\{0,1\}$ agree.

\smallskip\noindent\emph{Boolean connectives.}
The induction step for $\neg,\wedge,\vee,\to$ is immediate from the induction hypothesis.

\smallskip\noindent\emph{Modal case $\Box\psi$.}
Assume the induction hypothesis holds for $\psi$.
If $M_B,a\models \Box\psi$ for some $a\in B\setminus\{0,1\}$, then in particular $aRa$ and hence $M_B,a\models\psi$.
By the induction hypothesis, $\psi$ holds at every element of $B\setminus\{0,1\}$.
Also $aR1$ (Proposition~\ref{prop:1-related}), so $\Box\psi$ at $a$ implies $M_B,1\models\psi$.
Thus $\psi$ holds at every nonzero state, and therefore $\Box\psi$ holds at every $b\in B\setminus\{0,1\}$.
Conversely, if $\Box\psi$ fails at some $a\in B\setminus\{0,1\}$, then either $1$ or some nontrivial successor
witnesses failure of $\psi$, and the same reasoning shows $\Box\psi$ fails at all $b\in B\setminus\{0,1\}$.

\smallskip\noindent\emph{Modal case $\Diamond\psi$.}
If $M_B,a\models\Diamond\psi$ for some $a\in B\setminus\{0,1\}$, then there is $w$ with $aRw$ and $M_B,w\models\psi$.
Since $a\neq 0$, the witness $w$ cannot be $0$.
If $w\in B\setminus\{0,1\}$ then by induction $\psi$ holds at all nontrivial states, and hence every
$b\in B\setminus\{0,1\}$ satisfies $\Diamond\psi$ (witness $b$ itself).
If instead $w=1$, then $M_B,1\models\psi$; and since every $b\ne 0$ satisfies $bR1$, again every
$b\in B\setminus\{0,1\}$ satisfies $\Diamond\psi$.
The converse direction is similar.
\end{proof}

\begin{theorem}\label{thm:atomic-validities}
Let $B$ be a complete atomic Boolean algebra. In addition to the principles of
Theorem~\ref{thm:general-valid}, the following modal principles belong to $T_H(M_B)$:
\[
\Box\varphi \to \Box\Box\varphi\qquad (4),
\qquad
\Diamond\varphi \to \Box\Diamond\varphi\qquad (5),
\]
\[
\Diamond\Box\varphi \to (\varphi\to \Box\varphi)\qquad (W5).
%\qquad
%\Diamond\Box\varphi \to \Box\Diamond\varphi\qquad (.2).
\]
\end{theorem}

\begin{proof}
\emph{(4).}
Assume $M_B,a\models \Box\varphi$.
If $a=0$ there is nothing to show. Otherwise $a\ne 0$, hence $aRa$ and $aR1$.
Thus $M_B,a\models\varphi$ and $M_B,1\models\varphi$.
If $B=\Btwo$ then $R$ is trivially transitive and (4) holds.
If $B\ne\Btwo$, choose any $w\in B\setminus\{0,1\}$ with $aRw$ (take $w=a$ if $a\ne 1$, otherwise take any
nontrivial element). Then $M_B,w\models\varphi$, and by Theorem~\ref{thm:atomic-uniformity},
$\varphi$ holds at every element of $B\setminus\{0,1\}$.
Together with $M_B,1\models\varphi$, this shows that $\varphi$ holds at every nonzero state, hence every
nonzero state satisfies $\Box\varphi$. In particular, every successor of $a$ satisfies $\Box\varphi$, so
$M_B,a\models \Box\Box\varphi$.

\medskip\noindent
\emph{(5).}
Fix $a\in B$ and assume $M_B,a\models \Diamond\varphi$.
Then there exists $w\in B$ such that $aRw$ and $M_B,w\models \varphi$.

\smallskip\noindent\emph{Case 1: $w=1$.}
Let $b$ be any element with $aRb$.
By Proposition~\ref{prop:1-related} we have $bR1$, and since $M_B,1\models\varphi$, it follows that
$M_B,b\models\Diamond\varphi$ (witness $1$).  As $b$ was arbitrary, $M_B,a\models \Box\Diamond\varphi$.

\smallskip\noindent\emph{Case 2: $w\neq 1$.}
Then $w\in B\setminus\{0,1\}$.
By Theorem~\ref{thm:atomic-uniformity}, $\varphi$ holds at \emph{every} element of $B\setminus\{0,1\}$.
Let $b$ be any element with $aRb$.
If $b\neq 1$, then $b\in B\setminus\{0,1\}$, hence $M_B,b\models\varphi$, and since $bRb$ we get
$M_B,b\models\Diamond\varphi$.
If $b=1$, choose any $c\in B\setminus\{0,1\}$ (which exists when $B\ne \Btwo$). Then $1Rc$ and $M_B,c\models\varphi$,
so $M_B,1\models\Diamond\varphi$.
In either subcase $M_B,b\models\Diamond\varphi$, and since $b$ was arbitrary we obtain
$M_B,a\models\Box\Diamond\varphi$.

\medskip\noindent
\emph{(W5).}
Assume $M_B,a\models\Diamond\Box\varphi$ and $M_B,a\models\varphi$.
Choose $b$ with $aRb$ and $M_B,b\models \Box\varphi$.
Then $b\ne 0$, hence $bRb$ and $bR1$, so $M_B,b\models\varphi$ and $M_B,1\models\varphi$.
If $b\in B\setminus\{0,1\}$, then by Theorem~\ref{thm:atomic-uniformity} we have
$M_B,c\models\varphi$ for all $c\in B\setminus\{0,1\}$.
Together with $M_B,1\models\varphi$, this shows $\varphi$ holds at every nonzero state, and therefore
$M_B,a\models\Box\varphi$.
If instead $b=1$, then $M_B,1\models\Box\varphi$ already implies that $\varphi$ holds at every nonzero state,
so again $M_B,a\models\Box\varphi$.

%\medskip\noindent
%\emph{(.2).}
%Assume $M_B,a\models \Diamond\Box\varphi$ and choose $b$ with $aRb$ and $M_B,b\models\Box\varphi$.
%As above, $M_B,b\models\Box\varphi$ implies $M_B,1\models\varphi$ and $M_B,b\models\varphi$.
%If $b\in B\setminus\{0,1\}$ then by Theorem~\ref{thm:atomic-uniformity} $\varphi$ holds at all nontrivial states;
%together with $M_B,1\models\varphi$ this yields that every nonzero state $c$ satisfies $\Diamond\varphi$ (witness $c$).
%Since every successor of $a$ is nonzero, we obtain $M_B,a\models \Box\Diamond\varphi$.
%If $b=1$ then $M_B,1\models\Box\varphi$ already implies $\varphi$ holds at all nonzero states, and the same conclusion
%follows.
\end{proof}

\section{Completeness}\label{sec:completeness}

Section~\ref{sec:extended-semantics} was developed with respect to a fixed identification of propositional variables with parameter-free
set-theoretic sentences. For completeness, however, one should allow arbitrary assignments of
Boolean truth values to propositional variables. This leads naturally to a translation-based
semantics.

There is one further point to isolate. The Boolean value $0$ behaves degenerately: it has no
$R$-successors and it satisfies every atomic proposition, since $0\leq a$ for all $a\in B$. Thus the
present completeness theorem is obtained on the nonzero part of the algebra. In the terminology of Definition~\ref{def:theories} we are here interested in the notion of validity$^\sharp$. Accordingly, throughout
this section we work with
\[
B^+\ :=\ B\setminus\{0\}.
\]
On $B^+$ the co-consistency relation is reflexive as well as symmetric, and the resulting modal
logic is exactly $\KTB$.

\begin{definition}\label{def:B-translation-nonzero}
Fix a complete Boolean algebra $B$ and its Boolean-valued universe $V^{(B)}$.
\begin{enumerate}
    \item A \emph{$B$-translation} is a map $T:\Prop\to\Sent(L_B)$.
    \item Given $T$, define the induced valuation $v_T:\Prop\to B$ by
    \[
    v_T(p)\ :=\ \llbracket T(p)\rrbracket_B.
    \]
    \item If $W\subseteq B^+$ is nonempty, let $R$ be the inherited co-consistency relation on $W$:
    \[
    aRb \quad\Longleftrightarrow\quad a\wedge b\neq 0 \qquad (a,b\in W).
    \]
    We write $M_{B,W,T}=(W,R,v_T)$ for the corresponding Kripke model, where atomic satisfaction
    is given by
    \[
    M_{B,W,T},b\models p \quad\Longleftrightarrow\quad b\leq v_T(p).
    \]
\end{enumerate}
\end{definition}

If parameters are allowed, the Boolean truth-value map is surjective onto $B$.

\begin{lemma}\label{lem:truth-value-surjective}
For every $a\in B$ there exists an $L_B$-sentence $\sigma_a$ such that
\[
\llbracket \sigma_a\rrbracket_B=a.
\]
\end{lemma}

\begin{proof}
Fix $a\in B$. Let $u$ be the canonical $B$-name for the empty set, and let $v_a$ be the $B$-name
with $\mathrm{dom}(v_a)=\{u\}$ and $v_a(u)=a$. Then, by Definition~\ref{def:boolean-interpretation},
\[
\llbracket u\in v_a\rrbracket_B
=\bigvee_{x\in\mathrm{dom}(v_a)}\bigl(v_a(x)\wedge \llbracket x=u\rrbracket_B\bigr)
=v_a(u)\wedge 1
=a.
\]
Since $L_B$ contains constant symbols for $u$ and $v_a$, the sentence $\sigma_a$ may be taken to
be ``$u\in v_a$''.
\end{proof}

Let $\KTB$ denote the smallest normal modal logic extending $\mathsf{K}$ by the axioms
\[
(T)\quad \Box\varphi\to\varphi
\qquad\text{and}\qquad
(B)\quad \varphi\to\Box\Diamond\varphi.
\]
Equivalently, $\KTB$ is the logic of reflexive symmetric Kripke frames.

\begin{theorem}\label{thm:KTB-completeness}
For every modal formula $\varphi\in \Lbox$,
\[
\KTB\vdash \varphi
\quad\Longleftrightarrow\quad
M_{B,W,T}\models \varphi
\]
for all complete Boolean algebras $B$, all nonempty sets $W\subseteq B^+$, and all
$B$-translations $T$.
\end{theorem}


\begin{thebibliography}{99}
\bibitem[Bell]{Bell} J.~L.~Bell. \newblock \emph{Boolean-Valued Models and Independence Proofs in Set Theory}. \newblock Clarendon Press, Oxford, 1977.
\bibitem[Jech]{Jech} T.~Jech. \newblock \emph{Set Theory}. \newblock Springer Monographs in Mathematics. Springer, 2002.
\bibitem[LowePassmannTarafder]{LowePassmannTarafder} B.~L\"owe, R.~Pa{\ss}mann, and S.~Tarafder. \newblock Constructing illoyal algebra-valued models of set theory. \newblock \emph{Algebra Universalis}, 82(46), 2021.
\end{thebibliography}
\end{document}
